\section{Tau from Coupling via a Laplacian Generator}
\label{sec:tau-from-coupling}

\subsection{Pairwise coupling kernel}
Introduce a symmetric nonnegative pairwise kernel
\begin{equation}
A_{ij}=A_{ji}\ge 0,
\qquad A_{ii}=0.
\end{equation}
From this define
\begin{equation}
d_i = \sum_j A_{ij},
\qquad D = \operatorname{diag}(d_1,\dots,d_N),
\qquad L_A = D-A,
\end{equation}
which is the weighted graph Laplacian \cite{chung1997spectralgraph,fiedler1973algebraicconnectivity}.

\subsection{Logarithmic tau variables}
To guarantee positivity, write
\begin{equation}
\xi_i = \ln\tau_i.
\end{equation}
Impose the gauge-fixed linear system
\begin{equation}
L_A\xi = \kappa(d-\bar d\,\mathbf 1),
\qquad
\bar d = \frac1N\sum_i d_i,
\qquad
\sum_i \xi_i = 0.
\end{equation}

\subsection{Pseudoinverse solution}
Using the Moore--Penrose pseudoinverse \cite{penrose1955generalizedinverse}, the candidate solution is
\begin{equation}
\xi = \kappa L_A^+(d-\bar d\,\mathbf 1),
\end{equation}
and therefore
\begin{equation}
\tau_i = \tau_*\exp\!\Big(\kappa [L_A^+(d-\bar d\,\mathbf 1)]_i\Big).
\end{equation}

\subsection{Status}
This section records a first non-fit candidate generator of $\tau_i$ from pairwise coupling data. It is currently a registered hypothesis rather than a final derived law.
